% `tex.web` §890's widow penalty is an *argument* to `line_break`, and the two
% callers pass different ones.
%
%     procedure line_break(@!final_widow_penalty:integer);
%
% §1096's `end_graf` passes `\widowpenalty`. §1203 — the paragraph that a
% display interrupts — passes `\displaywidowpenalty`:
%
%     @<Go into display math mode@>=
%     begin if head=tail then {`\.{\\noindent\$\$}' or `\.{\$\${ }\$\$}'}
%       begin pop_nest; w:=-max_dimen;
%       end
%     else  begin line_break(display_widow_penalty);@/
%       ...
%
% So the line *before* a display carries `\interlinepenalty +
% \displaywidowpenalty`, and the second-to-last line of the paragraph resumed
% *after* it carries `\interlinepenalty + \widowpenalty`. Reading the parameter
% from one place gives both the same number, and nothing in plain or LaTeX's
% own defaults shows it, because they set the two alike.
%
% amsart does not: `\widowpenalty=10000` with `\displaywidowpenalty` left at
% LaTeX's 50. Every AMS paper with a display therefore has a page breakpoint of
% 50 where we offered 10000, and the page builder never considered it —
% `2608.17952v1` has eleven, and its running head landed 120 characters early
% because of them.
%
% The four numbers below are deliberately unequal, and no two of them sum to a
% third, so a wrong penalty cannot coincide with a right one:
%
%   14    = \interlinepenalty 3 + \clubpenalty 11      (first line)
%   3     = \interlinepenalty                          (ordinary)
%   53    = \interlinepenalty 3 + \displaywidowpenalty 50
%   10003 = \interlinepenalty 3 + \widowpenalty 10000
\input plain
\widowpenalty=10000 \displaywidowpenalty=50 \clubpenalty=11 \interlinepenalty=3
\showboxdepth=1 \showboxbreadth=40 \tracingonline=0
\setbox0=\vbox{\hsize=80pt
one two three four five six seven eight nine ten eleven twelve
$$x=1$$
after the display more words here to make another line or two of text\par}
\showbox0
\end
